\documentclass{article}

% Page size and margins
\usepackage[letterpaper,top=2cm,bottom=2cm,left=3cm,right=3cm,marginparwidth=1.75cm]{geometry}

% Useful packages
\usepackage{amsmath}
\usepackage{graphicx}
\usepackage[colorlinks=true, allcolors=blue]{hyperref}
\usepackage{lmodern}
\usepackage[T1]{fontenc}
\usepackage[english]{babel}

\title{Multi-Agent Reinforcement Learning for America's Cup Regattas}
\author{Your Team}
\date{\today}

\begin{document}
\maketitle

\begin{abstract}
In this project, we explore the application of Multi-Agent Reinforcement Learning (MARL)
to train sailboats to compete in a simulated America's Cup-style regatta. Using Proximal
Policy Optimization (PPO) in a vectorized multi-agent environment, we developed agents
capable of learning navigation strategies, managing foiling dynamics, and reacting to a
stochastic wind field. The final policy achieves a high race completion rate in evaluation,
while preserving realistic constraints such as gate crossings, boundary limits, and
collision-related penalties.
\end{abstract}

\section{Introduction}

This section will introduce the context of autonomous sailing and the motivation for using
Reinforcement Learning in competitive regatta scenarios. It will briefly describe the main
challenges of the task: stochastic wind, continuous boat control, tactical interaction between
boats, and the need to respect simplified racing rules.

\section{Background}

This section will present the theoretical foundations used in the project, including Markov
Decision Processes, policy optimization, and Proximal Policy Optimization. It will also
summarize the sailing concepts needed to understand the simulator, such as true wind angle,
velocity made good, polar performance curves, tacking, gybing, and foiling.

\section{Methodology}

This section will describe the system architecture and the formulation of the learning problem.
The discussion will include the simulated racecourse, the stochastic wind model, the boat
physics model, the multi-agent environment, the observation space, the continuous action
space, and the reward shaping strategy.

\section{Experimental Setup and Results}

This section will report the training configuration, PPO hyperparameters, evaluation protocol,
and empirical results. It will include quantitative metrics such as completion rate, reward
evolution, final distances, termination reasons, and qualitative observations from generated
race videos.

\section{Conclusions}

This section will summarize the achieved results and discuss the main limitations of the
current simulator. Possible future extensions include richer America's Cup rules, wind shadow
effects, more accurate hydrodynamics, independent policies for the two boats, and larger
multi-boat race scenarios.

\end{document}
